% GENERATED FILE. Edit canonical lessons, then run npm run generate:books.
% canonical-source-hash: fnv1a64:9cfd020fdbd8118b
% canonical-lessons: TE-C47-leg, TE-C47-tooth, TE-C47-hair, TE-C47-finger, TE-C47-stomach

\chapter{The Leg and the Tooth}
\label{ch:te-body-more}
\hlchaptermodality{47}

\begin{chapteropening}
\textbf{By the end of this chapter:} \emph{I can name five more parts of the body in Telugu, and say which of them held their old shape and which wore down past recognition.}

Complete TE-C47-stomach and demonstrate the chapter promise.
\end{chapteropening}

\section[kālu]{\te{కాలు} (kālu) — leg, foot}

\label{lesson:TE-C47-leg}



\begin{quote}
Before the new run: what did \emph{pustakaṁ} mean, and what did \emph{uppu} mean?
\end{quote}

\subsection*{You'll want to know: \te{కాలు}}
\textbf{\te{కాలు}} (\emph{kālu}) — \textquotedblleft{}leg, foot\textquotedblright{}.

Proto-Dravidian \textbf{*kāl}, and every sister kept it nearly whole: Tamil \ta{கால்}, Kannada \ka{ಕಾಲು}, Malayalam \ml{കാൽ}.

Set that beside \te{చెయ్యి}, the hand-word, which wore down so far that its family resemblance is hard to see at all. Two body-words of the same age; one held its shape and one did not, and nothing about the meaning predicts which.

\te{కాలు} covers the whole limb, foot and leg together. When precision is wanted, Sanskrit's \te{పాదం} names the foot alone.

\begin{grammarlens}[title={What you've built}]
The first of five more body words.
\end{grammarlens}

\subsection*{Your turn}
Say these aloud:

\begin{itemize}
\raggedright
  \item \emph{kālu}
  \item it once more, slowly
  \item \emph{kālu}, then \emph{tala}, so the top and the bottom sit together
\end{itemize}

\begin{tcolorbox}[breakable,colback=teal!4,colframe=teal!35!black,title={Before you move on}]
What does \te{కాలు} mean? (\textquotedblleft{}leg, foot\textquotedblright{}.) Say it once more.
\end{tcolorbox}
\section[pannu]{\te{పన్ను} (pannu) — a tooth}

\label{lesson:TE-C47-tooth}



\begin{quote}
Before the new one: what did \emph{kālu} mean?
\end{quote}

\subsection*{You'll want to know: \te{పన్ను}}
\textbf{\te{పన్ను}} (\emph{pannu}) — \textquotedblleft{}a tooth\textquotedblright{}.

Proto-Dravidian \textbf{*pal}. Tamil kept \ta{பல்}; Kannada has \ka{ಹಲ್ಲು}, with that \emph{p-} to \emph{h-} law working again exactly where the last chapter said to watch for it. Telugu took a third road, the \emph{-l-} coming out as \emph{-n-} and doubling.

An honest warning: \te{పన్ను} is also the word for a tax. Nothing joins the tooth to the tax but an accident of sound. Telugu has a fair number of these pairs, and context is the only fix — which is what context is for.

\begin{grammarlens}[title={What you've built}]
Two: a leg and a tooth.
\end{grammarlens}

\subsection*{Your turn}
Say these aloud:

\begin{itemize}
\raggedright
  \item \emph{pannu}
  \item it once more, slowly
  \item \emph{kālu}, then \emph{pannu}, and say which of the two also names something you pay
\end{itemize}

\begin{tcolorbox}[breakable,colback=teal!4,colframe=teal!35!black,title={Before you move on}]
What does \te{పన్ను} mean? (\textquotedblleft{}a tooth\textquotedblright{}.) Say it once more.
\end{tcolorbox}
\section[juṭṭu]{\te{జుట్టు} (juṭṭu) — hair (of the head)}

\label{lesson:TE-C47-hair}



\begin{quote}
Before the new one: what did \emph{pannu} mean?
\end{quote}

\subsection*{You'll want to know: \te{జుట్టు}}
\textbf{\te{జుట్టు}} (\emph{juṭṭu}) — \textquotedblleft{}hair (of the head)\textquotedblright{}.

\te{జుట్టు} is the hair on a head taken as one thing — the mass, not the strand. For a single hair Telugu reaches for a different noun, \te{వెంట్రుక}.

English leans on number for the same distinction --- \textquotedblleft{}a hair\textquotedblright{} against \textquotedblleft{}hair\textquotedblright{} --- while Telugu simply reaches for a different noun. \te{జుట్టు} also names the knot or tuft tied at the back of the head, which is most likely where the word started.

\begin{grammarlens}[title={What you've built}]
Three: a leg, a tooth, hair.
\end{grammarlens}

\subsection*{Your turn}
Say these aloud:

\begin{itemize}
\raggedright
  \item \emph{juṭṭu}
  \item it once more, slowly
  \item \emph{juṭṭu}, then \emph{tala}, and say which one sits on the other
\end{itemize}

\begin{tcolorbox}[breakable,colback=teal!4,colframe=teal!35!black,title={Before you move on}]
What does \te{జుట్టు} mean? (\textquotedblleft{}hair (of the head)\textquotedblright{}.) Say it once more.
\end{tcolorbox}
\section[vēlu]{\te{వేలు} (vēlu) — a finger}

\label{lesson:TE-C47-finger}



\begin{quote}
Before the new one: what did \emph{juṭṭu} mean?
\end{quote}

\subsection*{You'll want to know: \te{వేలు}}
\textbf{\te{వేలు}} (\emph{vēlu}) — \textquotedblleft{}a finger\textquotedblright{}.

Proto-Dravidian \textbf{*viral}. Tamil kept \ta{விரல்} entire; Kannada has \ka{ಬೆರಳು}; Telugu dropped the \emph{-r-} and came out with \te{వేలు}. On the page the three look barely related, and they are the same word.

It is still working inside compounds: \te{బొటనవేలు}, the thumb, is the stubby one of the set.

\begin{grammarlens}[title={What you've built}]
Four. Note that \te{చెయ్యి} and \te{వేలు} are separate roots — Telugu does not build the finger out of the hand.
\end{grammarlens}

\subsection*{Your turn}
Say these aloud:

\begin{itemize}
\raggedright
  \item \emph{vēlu}
  \item it once more, slowly
  \item \emph{vēlu}, then \emph{ceyyi}, and say which one has five of the other
\end{itemize}

\begin{tcolorbox}[breakable,colback=teal!4,colframe=teal!35!black,title={Before you move on}]
What does \te{వేలు} mean? (\textquotedblleft{}a finger\textquotedblright{}.) Say it once more.
\end{tcolorbox}
\section[kaḍupu]{\te{కడుపు} (kaḍupu) — the stomach, the belly}

\label{lesson:TE-C47-stomach}



\begin{quote}
Before the new one: what did \emph{vēlu} mean?
\end{quote}

\subsection*{You'll want to know: \te{కడుపు}}
\textbf{\te{కడుపు}} (\emph{kaḍupu}) — \textquotedblleft{}the stomach, the belly\textquotedblright{}.

\te{కడుపు} is the unfussy everyday word, and it carries the commonest way of saying you have eaten enough: \emph{kaḍupu niṇḍindi}, \textquotedblleft{}the stomach has filled\textquotedblright{}.

It also says, gently and without ever being explicit, that a woman is expecting. Sanskrit's \te{ఉదరం} exists in Telugu and stays where the clinical register lives.

\begin{grammarlens}[title={What you've built}]
Five more body words: a leg, a tooth, hair, a finger, a belly. With \te{తల}, \te{చెయ్యి}, \te{కన్ను}, \te{ముక్కు}, \te{నోరు} and \te{చెవి} behind them, you can point at most of a person.
\end{grammarlens}

\subsection*{Your turn}
Say these aloud:

\begin{itemize}
\raggedright
  \item \emph{kaḍupu}
  \item it once more, slowly
  \item all five in order, then add \emph{tala} and \emph{ceyyi} and say the seven as one list
\end{itemize}

\begin{tcolorbox}[breakable,colback=teal!4,colframe=teal!35!black,title={Before you move on}]
What does \te{కడుపు} mean? (\textquotedblleft{}the stomach, the belly\textquotedblright{}.) Say it once more.
\end{tcolorbox}
